\documentclass{article}
\usepackage{graphicx} % Required for inserting images
\usepackage{caption} % Required for \captionof
\usepackage{amsmath}
\usepackage{float}
\usepackage{comment}
\usepackage{microtype}


\title{Hőtani mérés}
\author{Kern Luca, Csongor Vincze}
\date{2026 Március 25}

\begin{document}

\maketitle

\section*{1. Feladat: A hőmérsékleti adatok leolvasása digitális multiméterről}

digitális multiméter: ellenállásmérési üzemmódba, megfelelő mérési tartományba állítottuk.
megfelelő kalibrációs meredekség:
mért ellenállás:
mért hőmérséklet:

\begin{comment}
\begin{table}[H]
    \centering
    \begin{tabular}{|c|c|}
        \hline
        Eltelt idő [s] & Hőmérséklet [°C] \\ \hline
    \end{tabular}
    \caption{Hőmérséklet az idő függvényében}
    \label{tab:hotan_1}
\end{table}
\end{comment}

\begin{figure}[H]
    \centering
    % \includegraphics[]{} % TODO: Add image path here
    \caption*{1. feladat: Mért hőmérséklet az idő függvényében}
    \label{}
\end{figure}

\section*{2. Feladat: A platinahőmérő kalibrálása}

Termosz színkódja:
Egyensúlyi hőmérséklet:
Soros ellenállás:
Hőmérsékleti korrekció:

\section*{3. Feladat: A termosz hőveszteségének mérése}

Meleg termosz:
Hideg termosz:
Hőmérsékletváltozás 1 perc alatt: $\Delta T=$

\section*{4. Feladat: A termosz hőkapacitásának becslése}

$m_{hideg}=$
$m_{meleg}=$
$c_{\text{rozsdamentes acél}}=500 \text{ J/(kg}^{\circ}\text{C)}$
termosz hőkapacitása: $C_{termosz}=$

\section*{5. Feladat: A víz fajhőjének mérése elektromos fűtéssel}

%fontos: előtte és utána is mérni 1-2 percig!!!

$m_{\text{víz}}=$
$T_{\text{egyensúlyi}}=$
$R_{\text{fűtőtest}}=$
$U_{\text{fűtőtest}}=$
$t=s$
termosz hőkapacitása az előző feladatból
$c_{\text{víz}}=$

%fajhő meghatározása egyenesillesztéssel

\section*{6. Feladat: Termosz hőkapacitásának mérése keveréssel}

$m_{hideg}=$
$m_{meleg}=$
$T_{hideg}=$
$T_{meleg}=$
$T_{\text{egyensúlyi}}=$
$c_{\text{víz}}= 4226 \text{ J/(kg}^{\circ}\text{C)}$   %irodalmi adat
$c_{\text{rozsdamentes acél}}=500 \text{ J/(kg}^{\circ}\text{C)}$
termosz hőkapacitása keveréssel: $C_{termosz}=$

\section*{7. Feladat Az alumínium fajhőjének mérése}

$d_{henger}=$
$h_{henger}=$
$\rho_{henger}=2,7 \text{ g/cm}^3$
$m_{\text{számolt}}=$
$m_{\text{mért}}=$
hideg termosz adatai:
$m_{\text{víz}}=$
$T_{\text{egyensúlyi}}=$
meleg termosz adatai:
$m_{\text{víz}}=$
$T_{\text{víz}}=$
$T_{\text{egyensúlyi}}=$

$c_{\text{víz}}= 4226 \text{ J/(kg}^{\circ}\text{C)}$   %irodalmi adat
$c_{\text{rozsdamentes acél}}=500 \text{ J/(kg}^{\circ}\text{C)}$
$C_{termosz}=$
$c_{Al}=$

\begin{figure}[H]
    \centering
    % \includegraphics[]{} % TODO: Add image path here
    \caption*{7. feladat: Mért hőmérséklet az idő függvényében}
    \label{}
\end{figure}

\section*{8. Feladat: A sárgaréz fajhőjének mérése}

$d_{henger}=$
$h_{henger}=$
$\rho_{henger}=2,7 \text{ g/cm}^3$
$m_{\text{számolt}}=$
$m_{\text{mért}}=$
hideg termosz adatai:
$m_{\text{víz}}=$
$T_{\text{egyensúlyi}}=$
meleg termosz adatai:
$m_{\text{víz}}=$
$T_{\text{víz}}=$
$T_{\text{egyensúlyi}}=$

$c_{\text{víz}}= 4226 \text{ J/(kg}^{\circ}\text{C)}$   %irodalmi adat
$c_{\text{rozsdamentes acél}}=500 \text{ J/(kg}^{\circ}\text{C)}$
$C_{termosz}=$
$c_{Cu}=$

\begin{figure}[H]
    \centering
    % \includegraphics[]{} % TODO: Add image path here
    \caption*{8. feladat: Mért hőmérséklet az idő függvényében}
    \label{}
\end{figure}


\end{document}